% Claim proof file for row claim_3_34 -- "ResumeLabelAlph".
% Auto-generated stub by scaffold/solving_current_row.py at row start. A
% manager agent dedicated to writing the TeX proof fills in the body below.
% The proof must:
%   - Stay close to the lecture notes' proof if one exists (always search
%     the LN first for a `\begin{proof}` block immediately following the
%     claim; copy / lightly edit when present).
%   - Otherwise use the LN paradigm and cite prior claims by their `ref`.
%   - Be self-contained enough that a mathematician can follow it without
%     re-reading the LN.
%   - Have no `sorry`, `omitted`, or "obvious" shortcuts.
%
% Compiles standalone via the `subfiles` package.

\documentclass[main]{subfiles}

\begin{document}
% ref: claim_3_34 (proof, with statement restated for readability)
% title: ResumeLabelAlph
\def\rowref{claim\_3\_34}
\def\rowtitle{ResumeLabelAlph}
\phantomsection\label{claim_3_34}
%
% Cross-references: see claim_<ref>_statement_<title>.tex in this folder
% for the corresponding statement.

% --- statement (restated) ---------------------------------------------------
\begin{Rem}
Let the assumptions be like in Theorem \refrow{def_3_23}.
    We also have the following rules: 
    \begin{enumerate}[resume,label=\alph*)]
        \setcounter{enumi}{12}
        \item Left Composition:
        \item[] $(A \iPerp  B \given  C) \land (D \iPerp  B \given C) \implies A \cup D \iPerp  B \given C.$
        \item Right Composition:
        \item[] $(A \iPerp  B \given  C) \land (A \iPerp  D \given C) \implies A \iPerp  B \cup  D \given C.$
         \item Left Intersection: If $A \cap D = \emptyset$ then:
        \item[] $(A \iPerp  B \given D \cup  C) \land (D \iPerp  B \given A \cup C) \implies A \cup D \iPerp  B \given C.$       
        \item Right Intersection: If $B \cap D = \emptyset$ then:
        \item[] $(A \iPerp  B \given D \cup  C) \land (A \iPerp  D \given B \cup C) \implies A \iPerp  B \cup D \given C.$
    \end{enumerate}
\end{Rem}

% --- proof ------------------------------------------------------------------
\begin{proof}
% TODO: write the proof body.
\end{proof}

\end{document}
